\documentclass{article}
\usepackage{amsmath, amssymb, amsthm}
\usepackage{hyperref}
\usepackage{geometry}
\geometry{margin=1in}

\title{Erd\H{o}s Problem \#769}
\date{Last updated: 2025-08-31}
\author{Source: erdosproblems.com}

\begin{document}
\maketitle

\noindent\textbf{Status:} OPEN \quad \textbf{No prize}

\noindent\textbf{Tags:} number theory, geometry

\medskip
\noindent\textbf{Problem Statement:}

Let $c(n)$ be minimal such that if $k\geq c(n)$ then the $n$-dimensional unit cube can be decomposed into $k$ homothetic $n$-dimensional cubes. Give good bounds for $c(n)$ - in particular, is it true that $c(n) \gg n^n$?

\bigskip
\noindent\textbf{Remarks:}

A problem first investigated by Hadwiger, who proved the lower bound\[c(n) \geq 2^n+2^{n-1}.\]It is easy to see that $c(2)=6$. Meier conjectured $c(3)=48$. Burgess and Erd\H{o}s \cite{Er74b} proved\[c(n) \ll n^{n+1}.\]Erd\H{o}s wrote 'I am certain that if $n+1$ is a prime then $c(n)>n^n$.'

Hudelson \cite{Hu98} proved that if $(2^n-1,3^n-1)=1$ then $c(n) < 6^n$, and in general $c(n) \ll (2n)^{n-1}$. Connor and Marmorino \cite{CoMa18} proved that\[c(n) \geq 2^{n+1}-1\]for all $n\geq 3$,\[c(n) \leq 1.8n^{n+1}\]if $n+1$ is prime, and\[c(n) \leq e^2n^n\]otherwise.

\bigskip
\noindent\textbf{References:}

[CoMa18] Connor, Peter and Marmorino, Phillip, Decomposing cubes into smaller cubes. J. Geom. (2018), Paper No. 19, 11.
 
 [Er74b] Erd\H{o}s, P., Remarks on some problems in number theory. Math. Balkanica (1974), 197-202.
 
 [Hu98] Hudelson, Matthew, Dissecting {$d$}-cubes into smaller {$d$}-cubes. J. Combin. Theory Ser. A (1998), 190--200.

\bigskip
\noindent\small{Source: \url{https://www.erdosproblems.com/769}}
\end{document}
